% The trust boundary as a table: what is proved vs. what is policy.
\begin{table}[t]
\centering
\renewcommand{\arraystretch}{1.25}
\begin{tabular}{@{}p{0.40\linewidth} l p{0.34\linewidth}@{}}
\toprule
\textbf{Concern} & \textbf{Status} & \textbf{Source of record} \\
\midrule
Group laws of $\mathbb{Z}^{\rdim}$            & proved   & \texttt{Vec.add\_comm/assoc/zero/neg} \\
Conservation is a subgroup                   & proved   & \texttt{conserves\_zero/add/neg} \\
Ledger fold preserves conservation           & proved   & \texttt{ledger\_conserves} \\
Double-entry (two-leg) balances              & proved   & \texttt{balanced\_pair} \\
Running \texttt{Money} arithmetic            & emitted  & \texttt{Ratio.Common.Emit} (from proofs) \\
\midrule
Normal-balance classification                & checked policy & \texttt{normalSide} (example, not invariant) \\
Posting rules / account mappings             & policy   & content-addressed control plane \\
Prices / FX / reference data                 & policy   & typed facts with provenance \\
NL intent $\to$ specification                 & untrusted & LLM proposal, compiled + checked \\
\bottomrule
\end{tabular}
\caption{The trust boundary. The top block is machine-checked (or emitted
from machine-checked sources) and forms the trusted computing base; the
bottom block is expressive policy that lives outside the kernel and cannot
violate conservation, because admission runs through
\cref{thm:conservation}.}
\label{tab:trust}
\end{table}
